% articulo-07-encabezado.tex - Encabezado y pie de página 
% (fancyhdr).
%
% Solo \section actualiza el encabezado (sin numeración).
% \subsection y niveles inferiores no tocan el encabezado.
%
% Variables YAML:
%   (sin declarar)                    →  verso: autor  /  recto: sección (por defecto)
%   encabezado-titulo-seccion: true   →  verso: título del documento  /  recto: sección
%   encabezado-seccion-seccion: true  →  verso y recto: título de sección
%   solo-folio: true                  →  solo número de página, sin texto
%   autor-recto: true                 →  con classoption: oneside, autor en el recto
%                                        (en lugar de sección; solo afecta al modo por defecto)
%
%   encabezado-versalitas: true  →  encabezados en versalitas
%   (sin declarar)               →  encabezados en cursiva (por defecto)
%
%   filete-encabezado: true      →  filete de 0.4 pt bajo el encabezado
%   (sin declarar)               →  sin filete (por defecto)
%
%   no-folio: true               →  páginas de inicio de sección sin folio
%   (sin declarar)               →  folio centrado en pie (por defecto, Chicago § 2.20)
%
%   titulos-sans: true           →  encabezados de página y folios en sans-serif,
%                                   coherente con la opción homónima de 13-secciones.tex.
%                                   Afecta a \headerfont y a los folios (\thepage) en
%                                   todos los estilos de página: fancy y plain
%                                   (páginas normales, preliminares y drop folio).
%   (sin declarar)               →  encabezados y folios en la familia principal (por defecto).
%
%   author: (metadato estándar)  →  nombre del autor (usado con autor-seccion)
%   title:  (metadato estándar)  →  título del documento (usado con encabezado-titulo-seccion)

\usepackage{fancyhdr}

% titlesec se carga aquí porque \sectionmark depende de él.
% El formato de \titleformat se configura en 13-secciones.tex.
\usepackage{titlesec}
\makeatletter
\def\sectionmark#1{\markright{#1}}
\def\subsectionmark#1{}
\def\subsubsectionmark#1{}
\makeatother

% \pagefont: familia tipográfica para encabezados y folios.
% Con titulos-sans: true, todo el aparato navegacional va en sans.
% Sin declarar, hereda la familia principal del documento.
$if(titulos-sans)$
\newcommand{\pagefont}[1]{{\sffamily#1}}
$else$
\newcommand{\pagefont}[1]{#1}
$endif$

% Estilo de texto para los encabezados.
% MakeLowercase evita que las mayúsculas queden más altas que las versalitas.
% \headerfont aplica primero \pagefont (familia) y luego el estilo (cursiva o versalitas).
$if(encabezado-versalitas)$
\newcommand{\headerfont}[1]{\pagefont{\textsc{\MakeLowercase{#1}}}}
$else$
\newcommand{\headerfont}[1]{\pagefont{\textit{#1}}}
$endif$

\pagestyle{fancy}
\fancyhf{}

% Folio: superior derecha en recto, superior izquierda en verso (Chicago § 1.6).
\fancyhead[RO]{\small\pagefont{\thepage}}
\fancyhead[LE]{\small\pagefont{\thepage}}

% Texto de encabezado según variables de contenido.
$if(solo-folio)$
% solo-folio: sin texto adicional.

$elseif(encabezado-titulo-seccion)$
% Verso: título del documento. Recto: título de sección.
$if(title)$
\fancyhead[RE]{\small\headerfont{$title$}}
$endif$
\fancyhead[LO]{\small\nouppercase{\headerfont{\rightmark}}}

$elseif(encabezado-seccion-seccion)$
% Verso y recto: título de sección.
\fancyhead[RE]{\small\nouppercase{\headerfont{\rightmark}}}
\fancyhead[LO]{\small\nouppercase{\headerfont{\rightmark}}}

$else$
% autor-seccion: comportamiento por defecto.
% Con autor-recto: true, el autor aparece en el recto (LO) en lugar de la sección.
% particularmente si se usa el classoption oneside, que solamente evalúa el recto
$if(author)$
$if(autor-recto)$
\fancyhead[RE]{\small\nouppercase{\headerfont{\rightmark}}}
\fancyhead[LO]{\small\headerfont{$for(author)$$author$$sep$, $endfor$}}
$else$
\fancyhead[RE]{\small\headerfont{$for(author)$$author$$sep$, $endfor$}}
\fancyhead[LO]{\small\nouppercase{\headerfont{\rightmark}}}
$endif$
$endif$
$endif$

% Filete bajo el encabezado.
$if(filete-encabezado)$
\renewcommand{\headrulewidth}{0.4pt}
$else$
\renewcommand{\headrulewidth}{0pt}
$endif$

% Páginas de inicio de sección (estilo plain).
% Drop folio centrado en pie por defecto (Chicago § 2.20).
% El encabezado está ausente; geometry reserva headheight+headsep
% igualmente, por lo que el cuerpo cae a ~2.54 cm del borde superior.
\fancypagestyle{plain}{%
  \fancyhf{}%
  $if(no-folio)$
  $else$
  \fancyfoot[C]{\small\pagefont{\thepage}}%
  $endif$
  \renewcommand{\headrulewidth}{0pt}%
}
